% Methodology chapter. Section 1 is the metric-ruler argument, moved from report 42
% (S1 metric integrity) and extended with the S4 bootstrap of the denominator (report 45).
% Written to be \input into the thesis body; it assumes a bibliography with the keys cited
% below (chapters/ref.bib carries stubs for them).

\chapter{Methodology}
\label{chap:methodology}

\section{Measuring accuracy on an intermittent series: the scaled-error ruler}
\label{sec:metric-ruler}

The three venues in this study trade on very different rhythms. The Beer Hall trades on
roughly five days in seven; Two River Taps on close to six; the village hall at Ellel on
about one and a fifth. Any accuracy figure that is to be compared \emph{across} those
venues, or against a backtest of the same venue, must therefore be scale-free. The Mean
Absolute Scaled Error (MASE) of \citet{hyndman_another_2006} is the natural choice: it
divides the forecast's mean absolute error by the in-sample mean absolute error of a naive
benchmark, so a value below one means the forecast beats that benchmark and the units of
revenue cancel.

The difficulty is that on a series with closed days the naive benchmark's denominator is
not unique, and the choices disagree by enough to reverse a conclusion. The seasonal-naive
benchmark differences the series at the seasonal lag; but ``the series'' can be read in more
than one way when a venue is shut for part of the week. This work makes the choice explicit
by naming four candidate bases and computing all four:

\begin{itemize}
  \item \texttt{calendar\_lag7}: the calendar-filled series (closed days entered as zero),
        differenced at lag seven.
  \item \texttt{trading\_lag7}: the trading-days-only series, differenced at lag seven.
  \item \texttt{trading\_same\_weekday}: the trading-days-only series differenced at the
        lag that lands on the same weekday once closures are compressed out.
  \item \texttt{calendar\_lag7\_active}: the calendar-filled series at lag seven, with any
        pair that has a closed endpoint dropped.
\end{itemize}

These are not cosmetic variants. On the Beer Hall the scale ranges from 316 to 632 across
the four; at Ellel it ranges from 180 to 806, a spread of four and a half to one. The same
forecast can be made to look four and a half times better or worse at Ellel by the choice of
denominator alone, and the earlier implementations recorded nowhere which choice had been
made.

The two natural implementations fail in opposite directions, which is why no single naive
repair would have sufficed. Compressing the closed days out (\texttt{trading\_lag7}) barely
shortens a series that already trades five days a week, so at the Beer Hall lag seven on the
trading index lands on the wrong weekday and roughly doubles the scale. Filling the closed
days with zeros (\texttt{calendar\_lag7}) has the opposite effect at Ellel: at one and a
fifth trading days a week, seventy-four per cent of the calendar lag-seven differences are
zero against zero, and the denominator collapses to about a quarter of its honest value.
The consequence for the headline is concrete. A live July forecast that had been published
at MASE $0.386$ was, on inspection, computed on \texttt{trading\_lag7}; on the reported
\texttt{calendar\_lag7} basis the same forecast scores $0.772$ against a backtest of
$0.745$, so the claim that it beat its own backtest by nearly half does not survive, and the
defensible statement is only that serving-horizon performance is consistent with the
backtest.

Two consequences follow for how a scaled metric must be reported. First, the basis must be
named and fixed: an unnamed MASE is uninterpretable, and comparing a forecast scored on one
basis with a backtest scored on another is the specific error above. This study fixes
\texttt{calendar\_lag7} as the reported basis, not because it is the most defensible of the
four (it is deflated by structural zeros) but because it is the only basis on which both the
rolling backtest and every out-of-sample confrontation already exist, so comparability
across the pre-registration chain is preserved. Second, and less obviously, a scaled metric
is a function of \emph{two} coordinates, not one: the denominator depends on the basis
\emph{and} on the training window it is estimated over. Because the store accumulates history,
the same basis yields a different scale as the window grows, so a published scaled figure is
reproducible only when both the basis and the training ceiling (an \texttt{as\_of} date) are
pinned. The confrontations in this work are therefore reported against a stated ceiling, and
the environment that produced them is version-pinned, so that a scaled number can be
regenerated exactly rather than approximately.

\subsection{How thin is the honest denominator? A bootstrap}
\label{sec:ruler-bootstrap}

The basis \texttt{calendar\_lag7\_active}, which drops closed-day pairs, is the most
defensible in principle: it is both weekday-aligned and undeflated. It is also the thinnest
in practice at exactly the venue whose intermittency motivates it. Of the lag-seven pairs it
keeps two hundred and seventy-six of three hundred and ninety-two at the Beer Hall, but only
twenty-eight of three hundred and eighty-five at Ellel. Whether a denominator estimated from
twenty-eight numbers can carry a gate is not a question argument can settle, so it is settled
by resampling.

For each venue and basis the absolute lag-difference vector, the numbers averaged into the
denominator, is resampled with replacement ten thousand times, and the scale's point
estimate, ninety-five per cent interval and interval width are recorded; that uncertainty is
then propagated into the venue's most recent confrontation MASE with the error numerator held
fixed. At the Beer Hall and Two River Taps every basis is estimated from a few hundred pairs
and the intervals are moderate, near thirty per cent of the point, so
\texttt{calendar\_lag7\_active} is adopted there: it is the cleanest basis and it is not, at
those venues, the noisiest. At Ellel the picture inverts. The active basis rests on twenty-
eight pairs and its ninety-five per cent scale interval spans plus or minus a third of the
point; the calendar basis, though wider-sampled, is deflated and induces a MASE whose own
interval runs from roughly $0.32$ to $0.55$ on scale uncertainty alone; and the trading bases
produce a denominator so large, because lag seven on the trading index reaches back almost
six weeks of highly variable demand, that the induced MASE falls near $0.09$ and reads as a
spurious near-perfect forecast. No basis is defensible at Ellel. At barely over one trading
day a week the seasonal-naive concept has no purchase, and the honest conclusion is that
scaled error is the wrong instrument there. This is precisely the case \citet{chatfield_all-zero_2007}
identify for highly intermittent series, where the lowest forecasting error need not give the
lowest system cost; at Ellel the reported quantity should be an unscaled error or a
cost-weighted one, not a MASE gate. The categorisation that motivates this treatment, and the
corrected constants used to apply it, are those of \citet{syntetos_categorization_2005} as
amended by \citet{kostenko_note_2006}.

% The diebold_comparing_1995 and harvey_testing_1997 keys cited below are expected in the
% Overleaf root ref.bib (standard forecasting-comparison references); add them there if absent.
\section{Choosing a model on a short series: fold count and the vanishing significance test}
\label{sec:fold-count}

Each venue is served by one forecaster, chosen from a ladder of nine by a rolling-origin
backtest. The original selection used six folds at a seven-day horizon. Six folds of seven
days is forty-two days of evidence, and forty-two days is not merely a thin sample: it is
exactly the sample size at which the standard test of whether one forecaster beats another
ceases to exist. The Diebold--Mariano test \citep{diebold_comparing_1995} compares two
forecasts through their mean loss differential; on a small sample the statistic is corrected
by the factor of \citet{harvey_testing_1997},
\[
  \mathrm{HLN}(n, h) = \sqrt{\frac{\,n + 1 - 2h + h(h-1)/n\,}{n}}.
\]
At a seven-day horizon the numerator $n + 1 - 14 + 42/n$ is zero at $n = 6$, so the
correction factor is exactly zero and the corrected statistic is undefined. This is an
algebraic identity, not a boundary approximation. The consequence is stark: the backtest that
selected every served model could not, even in principle, have attached a significance test to
the gap between two candidates. Selection without dispersion and selection on a forty-two-day
base are the same defect seen twice, because six disjoint seven-day windows are forty-two
days, forty-two days is $n = 6$, and $n = 6$ is where the test vanishes.

Widening the backtest from six disjoint windows to a step-one rolling origin, in which the
origin advances one day at a time, lifts the count to two hundred and seventy-three origins at
the Beer Hall, two hundred and sixty at Ellel, and two hundred and five at Two River Taps. At
those counts the correction factor recovers to about $0.97$ and a test becomes computable. The
gain is not free: consecutive step-one windows share six of their seven days, so their loss
differentials are serially correlated and are not independent draws. Any inference over them
therefore requires a block bootstrap or an autocorrelation-robust variance rather than a plain
$t$-test, which is the route the model confidence set of the next section takes.

\subsection{When the small sample chose the wrong model: the Beer Hall}
\label{sec:fold-count-beerhall}

Whether the fold count merely tightened intervals or actually reversed a conclusion is settled
at the Beer Hall, and the answer is the clearest single demonstration in this work of why the
sample size mattered. At the current store ceiling the six-fold window ranks the served
forecaster, \texttt{chronos2\_exo}, fifth of nine, and picks the day-of-week baseline
\texttt{robust\_dow} as the winner. That six-fold window falls on a high-variance summer
stretch of trading, a World Cup and a warm late June. Extend the same backtest to two hundred
and seventy-three origins and the served \texttt{chronos2\_exo} is restored to first. The small
sample gave the wrong answer and the large sample recovered the right one. Because the
mechanism is shown rather than asserted, one can say precisely how a forty-two-day evidence
base could mislead: it is a single draw of a noisy window, and the selection statistic, robbed
of its significance test by that very sample size, had no way to flag it as such.

A frozen control confirms the reading. Two River Taps closed on 8 May 2026, so its residual
frame did not grow as later months were ingested; its six-fold and full-origin rankings agree
to the digit and both confirm its served \texttt{ets} at every fold count. That isolates the
Beer Hall movement as a property of sample size on a live, growing series, not an artefact of
how the backtest was re-run. It also exposes a coupling that the metric-ruler section already
met from the other side: because ``the six most recent folds'' is a window anchored to the
store's last day, the same selection procedure yields a different table as the store grows, so
a published backtest figure is reproducible only against a stated training ceiling, exactly as
a scaled error is reproducible only against a stated basis. A backtest number, like a MASE,
carries two coordinates, and pinning only one of them is what let a forty-two-day window pass
for a settled result.

% The hansen_model_2011 key cited below is expected in the Overleaf root ref.bib (the model
% confidence set of Hansen, Lunde and Nason, Econometrica 2011); add it there if absent.
\section{Which forecasters the evidence cannot separate: the model confidence set}
\label{sec:mcs}

Recovering a computable test is not the same as declaring a winner. A rolling origin makes the
fold count large enough that a significance test exists, but the honest use of that test is not
to crown one forecaster; it is to report the set of forecasters the data cannot tell apart. The
instrument is the model confidence set of \citet{hansen_model_2011}, which returns, for a
stated confidence level, the collection of models whose predictive ability the evidence cannot
distinguish from the best. A single served model is defensible exactly when it lies inside its
own set, and the width of that set is itself the finding: a narrow set is a genuine ranking, a
wide one is an admission that the data does not support one.

Every choice the procedure needs was pre-registered before the set was computed, because a
confidence set has as many investigator degrees of freedom as any other test and fixing them
after seeing the answer would void it. The statistic is the range statistic $T_R$, the largest
standardised pairwise loss differential across the active models, with the matched elimination
rule that removes the model contributing that maximum. Significance is assessed by a moving
block bootstrap rather than an independent one, because consecutive step-one origins share six
of their seven days and their loss differentials are serially correlated by construction; an
independent resample would treat that overlap as fresh information and understate the variance.
The block length is set to seven, the horizon, which is the mechanical source of the overlap.
The bootstrap seed is fixed so the sets reproduce.

The confidence level is set at ten percent, and not at the conventional five, and this is a
deliberate choice about power rather than a relaxed standard. At these fold counts and with
top-of-ladder differences on the order of a few hundredths of a MASE, a five percent level has
no power to separate the leading models at all: it would return the full ladder at every venue
and say nothing. Ten percent is the tightest level at which the test can discriminate, and
naming it in advance is what keeps it from being the level chosen after the fact to produce a
congenial set. A secondary level of twenty-five percent is reported alongside; because a single
elimination sequence is thresholded at each level rather than recomputed, the twenty-five
percent set is a subset of the ten percent set by construction, so the pathological ordering in
which a model survives the stricter level but not the looser one is made impossible rather than
merely checked for.

\subsection{Why a test is possible at all: the paired variance}
\label{sec:mcs-paired}

Taken at face value the fold-level standard errors make the exercise look hopeless. The
per-fold loss of a single forecaster has a standard deviation several times larger than the
gaps between the leading forecasters, so any test of one mean loss against another, treating
the two as independent samples, would be powerless. The escape is that the forecasters are not
independent: they see the same demand on the same days, so their per-fold losses are strongly
correlated, and the standard deviation of their paired difference is far smaller than the
independent combination of their separate standard deviations would suggest. On this estate the
paired differential standard deviation runs from about a sixth to under a half of the
independent figure across the leading pairs, so the paired standard error is three to ten times
smaller than the marginal one. That gap is the entire reason a significance test is available
here, and it is why the bootstrap must resample folds jointly across models, applying one index
draw to every forecaster per replication: resampling each model independently would destroy the
cross-model correlation and inflate every set back to uselessness. The same correlation
structure justifies the block length empirically, since the autocorrelation of the paired
differential decays from a substantial lag-one value to about zero by the seventh lag, so a
block of seven is long enough to carry the dependence and no longer.

The result, applied across the served ladder, is that every served model lies inside its own
ninety percent set, and several of those sets are wide: the data separates the plainly weaker
rungs but cannot separate the foundation models from the served incumbent at the top. A wide
set is reported as the honest answer rather than being resolved by picking the numerically
smallest mean, because the range statistic's multiplicity correction is precisely what a naive
argmin lacks. Where the served choice cannot be shown superior on accuracy, the sentence worth
writing is that it rests on cold-start behaviour and inference cost rather than on a
demonstrated accuracy advantage; a spuriously precise winner would say less. The same
apparatus, reused unchanged, is what later adjudicates whether forecasting several venues
together under cross-series in-context learning improves on forecasting each alone, and returns
the same verdict in a different key: where the grouped and univariate arms cannot be
separated, the wider set is the finding, not a defect to be tuned away.
